%% LyX 2.4.4 created this file.  For more info, see https://www.lyx.org/.
%% Do not edit unless you really know what you are doing.
\documentclass[english]{article}
\usepackage[T1]{fontenc}
\usepackage[latin9]{inputenc}
\usepackage{enumitem}
\usepackage{amsmath}
\usepackage{amsthm}
\usepackage{amssymb}
\usepackage{cancel}
\usepackage{geometry}
\geometry{verbose,tmargin=1in,bmargin=1in,lmargin=1in,rmargin=1in}

\makeatletter
%%%%%%%%%%%%%%%%%%%%%%%%%%%%%% Textclass specific LaTeX commands.
\numberwithin{equation}{section}
\numberwithin{figure}{section}
\newlength{\lyxlabelwidth}      % auxiliary length 

%%%%%%%%%%%%%%%%%%%%%%%%%%%%%% User specified LaTeX commands.
\usepackage{fancyhdr}
\usepackage{lastpage}
\pagestyle{fancy}
\usepackage{enumitem}
\usepackage{ifthen}
%sets math fonts to be more readable
\everymath=\expandafter{\the\everymath\displaystyle}
%\DeclareMathSizes{10}{10}{10}{10}
\usepackage{multicol}

\usepackage{tikz}
\usetikzlibrary{plotmarks}
\usepackage{pgfplots}
\pgfplotsset{compat=newest}

\usepackage{calc}
\usepackage[nomessages]{fp}

\usepackage{advdate}    % Advancing/saving dates
\usepackage{datetime}   % Dates formatting
\usepackage{datenumber} % Counters for dates
% Added by lyx2lyx
\setlength{\parskip}{\smallskipamount}
\setlength{\parindent}{0pt}

\makeatother

\usepackage{babel}
\begin{document}
\renewcommand{\author}{Dr.\,James Ashe}

\newcommand{\classNum}{132}

\newcommand{\classSec}{C and K}

\newcommand{\nameType}{Name:}

\newcommand{\examType}{Worksheet 7.1 and 7.2}

\renewcommand{\date}{\formatdate{14}{4}{2014}}

%%First page's header%%%%%%%%%%%%%%%%%%%%%%
\fancypagestyle{firststyle} %%header settings for first pagr
{
	\fancyhead[L]{MTH \classNum{}(\classSec{}) \author{}\\
	\textbf{\examType{}} \date{}}
	\fancyhead[C]{\textbf{\nameType}}
	\setlength{\headsep}{14pt}
}
%%%%%%%%%%%%%%%%%%%%%%%%%%%%%%%%%%%%%%%%%%

%%Next page's header%%%%%%%%%%%%%%%%%%%%%%%
\setlist{nolistsep}
\lhead{\classNum{}(MTH \classSec{})} 
\chead{\textbf{\nameType}} 
\cfoot{\thepage{}~of~\pageref{LastPage}} 
\setlength{\headsep}{5pt} 
%%%%%%%%%%%%%%%%%%%%%%%%%%%%%%%%%%%%%%%%%%%

%%Various settings; see the preamble for other settings%%%%%
%%\setenumerate[0]{leftmargin=12pt,itemindent=0pt,labelwidth=0pt}
\renewcommand{\labelenumi}{\textbf{\arabic{enumi}.}}
\renewcommand{\labelenumii}{\textbf{\alph{enumii}.}}
\thispagestyle{firststyle}
%%%%%%%%%%%%%%%%%%%%%%%%%%%%%%%%%%%%%%%%%%
\newcommand{\xmax}{10}
\newcommand{\xmin}{-10}
\newcommand{\xstep}{0} %must be xmin+n
\newcommand{\ymax}{10}
\newcommand{\ymin}{-10}
\newcommand{\ystep}{0} %must be ymin+n

\newcommand{\xspace}{\vspace{.75cm}}

\subsection*{Sets and subsets}

\textbf{Write the following set in }\textbf{\emph{tabular}}\textbf{
and s}\textbf{\emph{et-builder}}\textbf{ form.}
\begin{enumerate}
\item The set of even, positive integers.\xspace\xspace
\end{enumerate}
\textbf{Decide whether the following are }\textbf{\emph{true}}\textbf{
or }\textbf{\emph{false}}\textbf{.}
\begin{enumerate}[resume]
\item $\ensuremath{\{1,2,3,4\}=\{4,3,1,2\}}$\xspace
\item $\ensuremath{\{a,b,\dots,z\}=\{b,c,\dots,z\}}$\xspace
\end{enumerate}
\textbf{Let $\ensuremath{A=\{2,4,6,10\}}$, $\ensuremath{B=\{2,10\}}$,
and $\ensuremath{C=\{2,4,6,10\}}$. Decide whether the following are
true or false.}
\begin{enumerate}[resume]
\item $\ensuremath{B\subset A}$\xspace
\item $\ensuremath{A\subseteq B}$\xspace
\item $\ensuremath{A\subset C}$\xspace
\item $\ensuremath{A\subseteq C}$ and $\ensuremath{C\subseteq A}$\xspace
\item $\ensuremath{A=C}$\xspace
\end{enumerate}
\textbf{Let $\ensuremath{E=\{2,4,6\}}$.}
\begin{enumerate}[resume]
\item How many subsets does $E$ have? List them.\xspace\xspace
\end{enumerate}

\subsection*{Set Operations and Venn Diagrams}

\textbf{Let $\ensuremath{A=\{2,4,6,10\}}$, $\ensuremath{B=\{2,4,8,10\}}$,
$\ensuremath{C=\{4,8,12\}}$, $\ensuremath{D=\{2,10\}}$, $\ensuremath{E=\{6\}}$,
and $\ensuremath{U=\{2,4,6,8,10,12,14\}}$. Insert $\ensuremath{\subseteq}$
or $\ensuremath{\nsubseteq}$ to make the statement true.}
\begin{enumerate}[resume]
\item $\ensuremath{A}\underline{\hspace{0.5cm}}\ensuremath{E}$\xspace
\item $\ensuremath{\emptyset}\underline{\hspace{0.5cm}}\ensuremath{A}$\xspace
\item $\ensuremath{A}\underline{\hspace{0.5cm}}\ensuremath{U}$\xspace
\item $\ensuremath{A}\underline{\hspace{0.5cm}}\ensuremath{U}$\xspace\newpage
\end{enumerate}
\textbf{Let $\ensuremath{U=\{1,2,3,4,5,6,7,8,9\}}$, $\ensuremath{X=\{2,4,6,8\}}$,
$\ensuremath{Y=\{2,3,4,5,6\}}$, and $\ensuremath{Z=\{1,2,3,8,9\}}$.
List the members of each set, using set braces.}
\begin{enumerate}[resume]
\item $\ensuremath{X\cup Y}$\xspace
\item $\ensuremath{X\cap Y}$\xspace
\item $X'$\xspace
\item $X'\cap Y'$\xspace\xspace
\item $X'\cap(Y'\cup Z)$\xspace\xspace
\item $(X\cap Y)\cup(X\cap Z)$\xspace\xspace
\item $(X\cap Y)\cap\{1,7\}$\xspace\xspace
\end{enumerate}
\textbf{Answer the following questions. }
\begin{enumerate}[resume]
\item There are $1500$ students at JCSU. If $100$ are math majors, $800$
are female, and $75$ are female math majors.

\begin{enumerate}[resume]
\item How many students are not math majors?\xspace
\item How many students are both math majors and female?\xspace\xspace
\item How many students are not math majors or female?\xspace\xspace
\end{enumerate}
\item Of the students in your class, 11 have sisters, 14 have brothers,
10 have a brother and a sister, and 5 have neither a brother or a
sister.

\begin{enumerate}
\item How many students have a brother or a sister?\xspace\xspace
\end{enumerate}
\begin{enumerate}[resume]
\item How many students are in the class? \xspace\xspace
\end{enumerate}
\end{enumerate}

\end{document}
